\section{Conclusions and future work}\label{sec:conclusion}
We built a complete, reproducible 5G SA system with three differentiated slices on a single
laptop, using only open-source software. Open5GS and UERANSIM ran natively on arm64 inside a
Lima VM on macOS, and the same scripts also target Windows through WSL2. Namespaces enforced
a realistic RAN/core/DN separation. The scenario,
architecture, traffic and measurements form one thread. The factory's three service classes
became three S-NSSAIs with their own DNN, pool, QoS profile and, for URLLC, their own UPF. The
signalling shows each UE receiving exactly its subscribed Allowed NSSAI and each PDU session
being anchored on the intended UPF. The KPIs show that the slice controls work and are
isolated. Capping the eMBB slice was enough to keep URLLC lossless with a p99 below 1.5\,ms during an
eMBB surge. Under deliberate platform overload, a dedicated URLLC UPF cut the share of
seconds that missed the 10\,ms p99 target from 65\,\% to 28\,\%, but it could not prevent losses in the shared RAN
components.
The sandbox exposes all of this to a non-expert in one page.

\paragraph{Future work.}
(1) Replace RLS with a PHY-level RAN (OpenAirInterface or srsRAN Project with its RF
simulator or an SDR) to study slice-aware MAC scheduling.
(2) Add an O-RAN near-RT RIC (FlexRIC or O-RAN SC) with an xApp that retunes slice resources
from the KPIs the sandbox already collects, closing the loop automatically.
(3) Enforce QoS by 5QI with PCC rules and GBR flows for URLLC.
(4) Move the NFs to Kubernetes/Helm with CPU pinning, so that the CPU artefacts of
\S\ref{sec:limits} can be controlled.
(5) Scale mIoT with a UE simulator that supports thousands of real registrations.
